\section{Benchmark: physics, realism and federation construction}
\label{sec:testbed}

No public dataset provides per-substation labelled fault and disturbance
waveforms suitable for a $K$-client federation; that gap is why the
benchmark generator exists. Rather than assert realism, every generator
parameter carries a verifiable source -- a normative standard, a public
dataset, or a peer-reviewed paper -- and is labelled by the \emph{kind}
of evidence behind it (\emph{measured-data}, \emph{standard},
\emph{fitted-synthetic}, \emph{modeling-choice}); the complete audit,
parameter by parameter, is Tables~\ref{tab:realism}--\ref{tab:realism2b}.
A test asserts that no row ships without a source, and
Section~\ref{sec:limitations} prices what the fitted-synthetic rows cost.

\subsection{Feeders, solver and events}

Five unmodified distribution test feeders are vendored -- the IEEE 13-,
34-, 37- and 123-bus feeders and the IEEE PES European Low Voltage Test
Feeder \cite{kersting2001} -- and clients are assigned feeders
cyclically, so feeder identity is a feature-skew axis before any label
skew. The IEEE 37-bus feeder is an ungrounded delta system, retained
deliberately: on it a ground fault appears as neutral displacement, not
fault current, so \emph{the same label means a different signal shape} at
different clients. The solver is the DSS-Extensions OpenDSS engine in
quasi-static daily mode at a \textbf{60\,s} step, metering at the
substation bay (feeder-head VT, first-section CT); a step that fails the
power flow is re-solved and then \emph{holds the last valid value} -- the
repeat-last convention of telemetry front ends -- at an aggregate held
rate of $0.0059\%$ (Table~\ref{tab:nonconv}).

Residential demand uses the CREST measured-model household profiles
\cite{richardson2010}; rooftop PV sits on $40\%$ of load buses at
$15\%$/$40\%$ penetration with Markov cloud transients fitted to measured
statistics \cite{hoke2013,woyte2007,lave2013}. Events arrive at a Poisson
rate per client-day -- line-to-ground, line-to-line, three-phase faults
at documented resistance distributions \cite{dagenhart2000,deandrade2010},
sags and swells targeted to land in the IEEE\,1159 bands
\cite{ieee1159,bollen2000}, and sustained cold-load-pickup overload
\cite{short2004} -- with episode durations of $2$--$10$ minutes.

\paragraph{No protection equipment is modelled}
The generator contains no relay, recloser, fuse or breaker: a fault stays
energized for its whole drawn duration where a real feeder would clear it
in cycles. Fault windows are therefore \emph{more} observable here than
in the field, absolute fault-class accuracies are optimistic, and the
benchmark supports the relative comparison across models and defenses,
not a field-accuracy claim.

\paragraph{The class prior is realistic, and that is the hard part}
The generator is tuned so $\ge 90\%$ of windows are \cls{normal}; the
realized share is $93.0$--$94.1\%$ across the battery federations
(Table~\ref{tab:dataset}). All primary metrics are computed under this
prior -- \macroF{} and per-class recall, never accuracy alone -- and the
rarest class can have single-digit support in a client split, so its
recall is unstable and reported, not smoothed.

\subsection{Measurement chain and federation construction}

The classifier sees instrumented telemetry, not clean solver output:
class-0.5 CT/VT error models \cite{iec61869-2,iec61869-3}, a 16-bit ADC,
dropout episodes at measured rates \cite{naspi2017}, and $60$\,s RMS
gateway telemetry -- not IEC 61850-9-2 sampled values, which the
generator does not model. System frequency is not a solver output, so
$f_{\mathrm{dev}}$ is a stationary Ornstein--Uhlenbeck process fitted to
measured statistics \cite{rydingorjao2020} with event-coupled dips -- a
fitted-synthetic channel whose possible label-shortcut is audited in
Section~\ref{sec:res-frame}.

Windows are $T{=}64$ samples at stride~16 ($75\%$ overlap), $87$ per
simulated day, $2\,610$ per client at the 30-day default. Heterogeneity
is two-axis: feeder/load-shape identity gives feature skew and a
Dirichlet$(\alpha)$ draw over the eight classes gives label skew
\cite{hsu2019} ($\alpha=0.5$ reference, $\{0.1,1.0\}$ and a near-uniform
control swept). Splits are per-client $70/15/15$, seeded, the normalizer
fitted on the training split only. Two protocols exist and both are
reported: the pre-registered seeded \emph{random} permutation -- under
which the committed \texttt{experiments/measure\_leakage.py} finds
$99.9\%$ of test windows overlapping a training window -- and the
leakage-controlled \emph{chronological} split (contiguous blocks,
four-window purge, $\approx 0\%$ overlap), adopted as the headline
protocol. The consensus layer's latency is a sourced wide-area lognormal
(median 25\,ms) \cite{aikat2003}. The same pipeline runs
unmodified on the NREL SMART-DS systems \cite{smartds}, the declared
scale-up path.
